% Kanonische Quelle: Band 19 und Beweistabellen verwenden denselben Text.
\chapter{Dedekindsche Schnitte}

\section{Schnittprädikat und Trägermenge}

\FormulaDefDeltaK[Dedekindscher Schnitt der rationalen Zahlen]{%
  \RealCut{A}%
}{DedekindCutPredicate}{%
  \DeltaRow{Mengen}{A\dsep p\dsep q\dsep r}
  \DeltaRow{Relationsoperatoren}{<}
  \DeltaPrem{Geordneter rationaler Körper}{\mathbb Q}
  \DeltaRow{\textbf{Axiome}}{}
  \DeltaPrem{Rationale Teilmenge}{A\subseteq\mathbb Q}
  \DeltaPrem{Nichtleerheit}{A\neq\varnothing}
  \DeltaPrem{Echtheit}{A\neq\mathbb Q}
  \DeltaRow{Nach unten abgeschlossen}{%
    p\in A\dsep q\in\mathbb Q\dsep q<p\vdash q\in A}
  \DeltaRow{Kein größtes Element}{p\in A\vdash\exists r\in A\,(p<r)}
  \DeltaRow{\textbf{Neues Symbol}}{}
  \DeltaPrem{Schnittprädikat}{\RealCut{A}}
}

\FormulaAxiomDeltaK[Rationale Trägermenge]{%
  \RealCut{A}\vdash A\subseteq\mathbb Q%
}{DedekindCutSubsetAxiom}{%
  \DeltaRow{Mengen}{A\dsep p\dsep q\dsep r}
}

\FormulaAxiomDeltaK[Nichtleerheit]{%
  \RealCut{A}\vdash A\neq\varnothing%
}{DedekindCutNonemptyAxiom}{%
  \DeltaRow{Mengen}{A\dsep p\dsep q\dsep r}
}

\FormulaAxiomDeltaK[Echter Schnitt]{%
  \RealCut{A}\vdash A\neq\mathbb Q%
}{DedekindCutProperAxiom}{%
  \DeltaRow{Mengen}{A\dsep p\dsep q\dsep r}
}

\FormulaAxiomDeltaK[Abwärtsabgeschlossenheit]{%
  \RealCut{A}\vdash \forall p\in A\,\forall q\in\mathbb Q\,(q<p\rightarrow q\in A)%
}{DedekindCutDownwardAxiom}{%
  \DeltaRow{Mengen}{A\dsep p\dsep q\dsep r}
}

\FormulaAxiomDeltaK[Kein größtes Element]{%
  \RealCut{A}\vdash \forall p\in A\,\exists r\in A\,(p<r)%
}{DedekindCutNoMaximumAxiom}{%
  \DeltaRow{Mengen}{A\dsep p\dsep q\dsep r}
}

\FormulaThmDeltaK[Schnittkriterium]{%
  \begin{aligned}[t]
  \RealCut{A}\eqvdash{}&
    A\subseteq\mathbb Q\land A\neq\varnothing\land A\neq\mathbb Q
    \land{}\\[-2pt]
  &\forall p\in A\,\forall q\in\mathbb Q\,
      \bigl(q<p\rightarrow q\in A\bigr)
    \land{}\\[-2pt]
  &\forall p\in A\,\exists r\in A\,(p<r).
  \end{aligned}%
}{DedekindCutCharacterization}{%
  \DeltaRow{Mengen}{A\dsep p\dsep q\dsep r}
  \DeltaRow{Relationsoperatoren}{<}
}
\begin{tabproofsplitwide}
  \proofpartwide{Hinrichtung}
    \proofstepwidestar[1]{\RealCut{A}}{\rA}
    \proofstepwidestar[1]{A\subseteq\mathbb Q}{%
      \FormulaRefAuto{DedekindCutSubsetAxiom}[ax]{1}}
    \proofstepwidestar[1]{A\neq\varnothing}{%
      \FormulaRefAuto{DedekindCutNonemptyAxiom}[ax]{1}}
    \proofstepwidestar[1]{A\neq\mathbb Q}{%
      \FormulaRefAuto{DedekindCutProperAxiom}[ax]{1}}
    \proofstepwidestar[1]{\forall p\in A\,\forall q\in\mathbb Q\,(q<p\rightarrow q\in A)}{%
      \FormulaRefAuto{DedekindCutDownwardAxiom}[ax]{1}}
    \proofstepwidestar[1]{\forall p\in A\,\exists r\in A\,(p<r)}{%
      \FormulaRefAuto{DedekindCutNoMaximumAxiom}[ax]{1}}
    \proofstepwidestar[1]{%
      \begin{aligned}[t]
      &A\subseteq\mathbb Q\land A\neq\varnothing\land A\neq\mathbb Q\land{}\\[-2pt]
      &\forall p\in A\,\forall q\in\mathbb Q\,(q<p\rightarrow q\in A)\land{}\\[-2pt]
      &\forall p\in A\,\exists r\in A\,(p<r)
    \end{aligned}}{\rAIStar{2,3,4,5,6}}
  \closeproofpartwide

  \proofpartwide{Rückrichtung}
    \proofstepwidestar[1]{%
      \begin{aligned}[t]
      &A\subseteq\mathbb Q\land A\neq\varnothing\land A\neq\mathbb Q\land{}\\[-2pt]
      &\forall p\in A\,\forall q\in\mathbb Q\,(q<p\rightarrow q\in A)\land{}\\[-2pt]
      &\forall p\in A\,\exists r\in A\,(p<r)
    \end{aligned}}{\rA}
    \proofstepwidestar[1]{A\subseteq\mathbb Q}{\rAEa{1}}
    \proofstepwidestar[1]{A\neq\varnothing}{\rAEa{\rAEb{1}}}
    \proofstepwidestar[1]{A\neq\mathbb Q}{\rAEa{\rAEb{\rAEb{1}}}}
    \proofstepwidestar[1]{\forall p\in A\,\forall q\in\mathbb Q\,(q<p\rightarrow q\in A)}{\rAEa{\rAEb{\rAEb{\rAEb{1}}}}}
    \proofstepwidestar[1]{\forall p\in A\,\exists r\in A\,(p<r)}{\rAEb{\rAEb{\rAEb{\rAEb{1}}}}}
    \proofstepwidestar[1]{\RealCut{A}}{%
      \FormulaRefAuto{DedekindCutPredicate}[def]{2,3,4,5,6}}
  \closeproofpartwide
\end{tabproofsplitwide}

Neben der rationalen Typisierung sind die vier Schnittbedingungen wesentlich. Die
Abwärtsabgeschlossenheit macht einen Schnitt zu einem rationalen
Anfangsabschnitt; die letzte Bedingung entfernt die ansonsten mögliche
Mehrdeutigkeit am Rand eines rationalen Anfangsabschnitts.

\FormulaDefDeltaK[Die Menge der reellen Zahlen]{%
  \mathbb R\coloneqq
  \bigl\{A\in\powerset(\mathbb Q)\bigm|\RealCut{A}\bigr\}%
}{RealCutCarrier}{%
  \DeltaRow{Mengen}{A\dsep\mathbb R}%
  \DeltaRow{Prädikate}{\RealCut{\,\cdot\,}}%
  \DeltaRow{\textbf{Neues Symbol}}{}%
  \DeltaRow{Reelle Zahlen}{\mathbb R}%
}

\FormulaThmDeltaK[Elementcharakterisierung der reellen Zahlen]{%
  A\in\mathbb R\eqvdash\RealCut{A}%
}{RealCutMembership}{%
  \DeltaRow{Mengen}{A}%
}
\begin{tabproofsplitwide}
  \proofpartwideindR[Hinrichtung]{%
    A\in\mathbb R\rightarrow\RealCut A%
  }{A\in\mathbb R\rightarrow\RealCut A}[RealCutMembershipForward]
    \proofstepwidestar[1]{A\in\mathbb R}{\rA}
    \proofstepwidestar[]{%
      \mathbb R=\{X\in\powerset(\mathbb Q)\mid\RealCut X\}}{%
      \FormulaRefAuto{RealCutCarrier}[def]}
    \proofstepwidestar[1]{%
      A\in\powerset(\mathbb Q)\land\RealCut A}{%
      \rIE{1,2}}
    \proofstepwidestar[1]{\RealCut A}{\rAEb{3}}
  \proofstepwidestar[]{A\in\mathbb R\rightarrow\RealCut A}{\rRI{1,4}}
  \closeproofpartwideindR

  \proofpartwideindR[Rückrichtung]{%
    \RealCut A\rightarrow A\in\mathbb R%
  }{\RealCut A\rightarrow A\in\mathbb R}[RealCutMembershipBackward]
    \proofstepwidestar[1]{\RealCut A}{\rA}
    \proofstepwidestar[1]{A\subseteq\mathbb Q}{%
      \FormulaRefAuto{DedekindCutCharacterization}[thm]{1}}
    \proofstepwidestar[1]{A\in\powerset(\mathbb Q)}{%
      \FormulaRefAuto{x\in\powerset(A)\eqvdash x\subseteq A}{2}}
    \proofstepwidestar[1]{%
      A\in\powerset(\mathbb Q)\land\RealCut A}{\rAI{3,1}}
    \proofstepwidestar[1]{A\in\mathbb R}{%
      \FormulaRefAuto{RealCutCarrier}[def]{4}}
  \proofstepwidestar[]{\RealCut A\rightarrow A\in\mathbb R}{\rRI{1,5}}
  \closeproofpartwideindR

  \proofpartwide{Zusammenfassung}
    \proofstepwidestar[]{A\in\mathbb R\rightarrow\RealCut A}{%
      \FormulaRefAuto{RealCutMembershipForward}[thm]}
    \proofstepwidestar[]{\RealCut A\rightarrow A\in\mathbb R}{%
      \FormulaRefAuto{RealCutMembershipBackward}[thm]}
    \proofstepwidestar[]{A\in\mathbb R\leftrightarrow\RealCut A}{%
      \rLRI{1,2}}
  \closeproofpartwide
\end{tabproofsplitwide}

Diese Charakterisierung erlaubt es im Folgenden, eine reelle Zahl und den sie
darstellenden Schnitt ohne zusätzliche Klassenklammern miteinander zu
identifizieren.

\section{Hauptschnitte und rationale Einbettung}

\FormulaDefDeltaK[Hauptschnitt einer rationalen Zahl]{%
  q\in\mathbb Q\vdash
  \RealHat q\coloneqq\{p\in\mathbb Q\mid p<q\}%
}{RealPrincipalCut}{%
  \DeltaRow{Mengen}{p\dsep q}%
  \DeltaRow{Relationsoperatoren}{<}%
  \DeltaRow{\textbf{Neues Symbol}}{}%
  \DeltaRow{Hauptschnitt}{\RealHat q}%
}

\FormulaThmDeltaK[Hauptschnitte sind Dedekindsche Schnitte]{%
  q\in\mathbb Q\vdash\RealCut{\RealHat q}%
}{RealPrincipalCutIsCut}{%
  \DeltaRow{Mengen}{p\dsep q\dsep r}%
}
\begin{tabproofsplitwide}
\proofpartwide{Nichtleerheit}
\proofstepwidestarbreakkeep[1]{q\in\mathbb Q}{\rA}
    \proofstepwidestarbreakkeep[]{\IntOne\in\PositiveIntegers}{%
      \FormulaRefAuto{IntegerOnePositive}[thm]}
    \proofstepwidestarbreakkeep[]{\IntZero<\IntOne}{%
      \FormulaRefAuto{PositiveIntegersDef}[def]{2}}
    \proofstepwidestarbreakkeep[]{%
      \IntZero\leq\IntOne\land\IntZero\neq\IntOne}{%
      \FormulaRefAuto{IntegerStrictOrderDef}[def]{3}}
    \proofstepwidestarbreakkeep[]{\IntZero\in\mathbb Z}{%
      \FormulaRefAuto{IntZeroCarrier}[thm]}
    \proofstepwidestarbreakkeep[]{\IntOne\in\mathbb Z}{%
      \FormulaRefAuto{IntOneCarrier}[thm]}
    \proofstepwidestarbreakkeep[]{%
      \RatEmbed(\IntZero)\leq\RatEmbed(\IntOne)}{%
      \FormulaRefAuto{IntegerRationalEmbeddingOrder}[thm]{4,5,6}}
    \proofstepwidestarbreakkeep[]{%
      \RatEmbed(\IntZero)\neq\RatEmbed(\IntOne)}{%
      \FormulaRefAuto{IntegerRationalEmbeddingInjective}[thm]{4,5,6}}
    \proofstepwidestar[]{\RatOne=\RatEmbed(\IntOne)}{%
      \FormulaRefAuto{RationalOneDef}[def]}
    \proofstepwidestar[]{\RatZero=\RatEmbed(\IntZero)}{%
      \FormulaRefAuto{RationalZeroDef}[def]}
    \proofstepwidestarbreakkeep[]{%
      \RatZero=\RatEmbed(\IntZero)\land
      \RatOne=\RatEmbed(\IntOne)}{%
      \rAI{10,
        9}}
    \proofstepwidestarbreakkeep[]{\RatZero<\RatOne}{%
      \FormulaRefAuto{RationalStrictOrderDef}[def]{7,8,11}}
    \proofstepwide[1,12]{q-\RatOne}{<}{q}{%
      \FormulaRefAuto{RationalOrderTranslation}[thm]{1,12}}
    \proofstepwidestarbreakkeep[1,12]{%
      \exists p\in\mathbb Q\,(q-\RatOne<p\land p<q)}{%
      \FormulaRefAuto{RationalOrderDense}[thm]{13}}
    \proofstepwidestarbreakkeep[1,12]{p\in\mathbb Q\land p<q}{\rA}
    \proofstepwidestarbreakkeep[1,12]{p\in\RealHat q}{%
      \FormulaRefAuto{RealPrincipalCut}[def]{15}}
    \proofstepwidestarbreakkeep[1,12]{\RealHat q\neq\varnothing}{%
      \FormulaRefAuto{x\in A\vdash A\neq\varnothing}{16}}

\closeproofpartwide

\proofpartwide{Echtheit}
  \setcounter{proofstepnr}{17}% Fortlaufende Schritte dieses Beweises.
\proofstepwidestar[1]{\TotOrd{\mathbb Q,\leq}}{%
      \FormulaRefAuto{RationalOrderTotal}[thm]{1}}
\proofstepwidestarbreakkeep[1]{q\notin\RealHat q}{%
      \FormulaRefAuto{RealPrincipalCut}[def]{%
        18}}
    \proofstepwidestarbreakkeep[1]{\RealHat q\neq\mathbb Q}{%
      \FormulaRefAuto{x\in B, x\not\in A\vdash A\neq B}{1,19}}

\closeproofpartwide

\proofpartwideindR[Nichtleerheit und Echtheit]{%
    q\in\mathbb Q\vdash
    \RealHat q\neq\varnothing\land\RealHat q\neq\mathbb Q
  }{%
    q\in\mathbb Q\vdash
    \RealHat q\neq\varnothing\land\RealHat q\neq\mathbb Q
  }[RealPrincipalCutNonemptyProper]
  \setcounter{proofstepnr}{20}% Fortlaufende Schritte dieses Beweises.
\proofstepwidestarbreakkeep[1]{%
      \RealHat q\neq\varnothing\land\RealHat q\neq\mathbb Q}{%
      \rAI{17,20}}
    \proofstepwidestarbreakkeep[1]{%
      \RealHat q\neq\varnothing\land\RealHat q\neq\mathbb Q}{%
      \rEE{14,15,21}}

\closeproofpartwideindR


  \proofpartwideindR[Abwärtsabschluss]{%
    \begin{aligned}[t]
      &q\in\mathbb Q\dsep p\in\RealHat q\dsep s\in\mathbb Q\dsep s<p\\[-2pt]
      &\quad\vdash s\in\RealHat q
    \end{aligned}
  }{%
    q\in\mathbb Q\dsep p\in\RealHat q\dsep s\in\mathbb Q\dsep s<p
    \vdash s\in\RealHat q
  }[RealPrincipalCutDownwardClosed]
    \proofstepwidestar[1]{q\in\mathbb Q}{\rA}
    \proofstepwidestar[2]{p\in\RealHat q}{\rA}
    \proofstepwidestar[3]{s\in\mathbb Q}{\rA}
    \proofstepwidestar[4]{s<p}{\rA}
    \proofstepwide[1,2]{p}{<}{q}{%
      \FormulaRefAuto{RealPrincipalCut}[def]{1,2}}
    \proofstepwide[1,2,4]{s}{<}{q}{%
      \FormulaRefAuto{RationalOrderTotal}[thm]{4,5}}
    \proofstepwidestar[1,2,3,4]{s\in\RealHat q}{%
      \FormulaRefAuto{RealPrincipalCut}[def]{3,6}}
  \closeproofpartwideindR

  \proofpartwideindR[Kein größtes Element]{%
    q\in\mathbb Q\dsep p\in\RealHat q
    \vdash\exists r\in\RealHat q\,(p<r)
  }{%
    q\in\mathbb Q\dsep p\in\RealHat q
    \vdash\exists r\in\RealHat q\,(p<r)
  }[RealPrincipalCutNoMaximum]
    \proofstepwidestar[1]{q\in\mathbb Q}{\rA}
    \proofstepwidestar[2]{p\in\RealHat q}{\rA}
    \proofstepwide[1,2]{p}{<}{q}{%
      \FormulaRefAuto{RealPrincipalCut}[def]{1,2}}
    \proofstepwidestar[1,2]{%
      \exists r\in\mathbb Q\,(p<r\land r<q)}{%
      \FormulaRefAuto{RationalOrderDense}[thm]{3}}
    \proofstepwidestar[1,2]{r\in\mathbb Q\land p<r\land r<q}{\rA}
    \proofstepwidestar[1,2]{r\in\RealHat q}{%
      \FormulaRefAuto{RealPrincipalCut}[def]{5}}
    \proofstepwidestar[1,2]{\exists r\in\RealHat q\,(p<r)}{%
      \rEI{5,6}}
    \proofstepwidestar[1,2]{\exists r\in\RealHat q\,(p<r)}{%
      \rEE{4,5,7}}
  \closeproofpartwideindR

  \proofpartwide{Schluss}
    \proofstepwidestar[1]{q\in\mathbb Q}{\rA}
    \proofstepwidestar[1]{\RealHat q\subseteq\mathbb Q}{%
      \FormulaRefAuto{RealPrincipalCut}[def]{1}}
    \proofstepwidestar[1]{%
      \RealHat q\neq\varnothing\land\RealHat q\neq\mathbb Q}{%
      \FormulaRefAuto{RealPrincipalCutNonemptyProper}[thm]{1}}
    \proofstepwidestar[1]{%
      \forall p\in\RealHat q\,\forall s\in\mathbb Q\,
      (s<p\rightarrow s\in\RealHat q)}{%
      \FormulaRefAuto{RealPrincipalCutDownwardClosed}[thm]{1}}
    \proofstepwidestar[1]{%
      \forall p\in\RealHat q\,\exists r\in\RealHat q\,(p<r)}{%
      \FormulaRefAuto{RealPrincipalCutNoMaximum}[thm]{1}}
    \proofstepwidestar[1]{\RealCut{\RealHat q}}{%
      \FormulaRefAuto{DedekindCutCharacterization}[thm]{2,3,4,5}}
  \closeproofpartwide
\end{tabproofsplitwide}

\FormulaDefDeltaK[Kanonische Einbettung der rationalen Zahlen]{%
  \begin{aligned}[t]
    &\RealEmb\colon\mathbb Q\to\mathbb R,\\[-2pt]
    &\forall q\in\mathbb Q\,
      \bigl(\RealEmb(q)\coloneqq\RealHat q\bigr)
  \end{aligned}%
}{RationalToRealEmbedding}{%
  \DeltaRow{Mengen}{p\dsep q}%
  \DeltaRow{Funktionen}{\RealEmb}%
  \DeltaRow{\textbf{Neues Symbol}}{}%
  \DeltaRow{Rationale Einbettung}{\RealEmb}%
}
\begin{tabproof}
  \proofstep{1}{q\in\mathbb Q}{\rA}
  \proofstep{1}{\RealCut{\RealHat q}}{%
    \FormulaRefAuto{RealPrincipalCutIsCut}[thm]{1}}
  \proofstep{1}{\RealHat q\in\mathbb R}{%
    \FormulaRefAuto{A\in\mathbb R\eqvdash\RealCut{A}}[thm]{2}}
  \proofstep{}{\forall q\in\mathbb Q\,\bigl(\RealHat q\in\mathbb R\bigr)}{%
    \rUI{\rRI{1,3}}}
\end{tabproof}

\FormulaThmDeltaK[Die rationale Einbettung ist eine Funktion]{%
  \RealEmb\colon\mathbb Q\to\mathbb R%
}{RationalToRealEmbeddingFunction}{%
  \DeltaRow{Funktionen}{\RealEmb}%
}
\begin{tabproof}
  \proofstep{}{\RealEmb\colon\mathbb Q\to\mathbb R}{%
    \FormulaRefAuto{RationalToRealEmbedding}[def]}
\end{tabproof}

